\begin{itemize}
\item \textbf{Calidad: ACO es mejor en todos los casos medidos.} Con 2\,048 agentes ACO obtiene menor $L_{mejor}$ en 15 de 15 combinaciones de $n$ e iteraciones, y en las 36 configuraciones emparejadas (mismo $n$, agentes e iteraciones; $n\le20\,000$) gana en 36 de 36. En $n=20$ ACO llega al óptimo de Held-Karp ($L/L_{ref}=1,000$ con 2\,048 agentes), mientras PSO queda 5,0\,\% por encima con 100 it. y 54,8\,\% con 5 it. En los demás $n$ ACO queda como máximo 3,4\,\% sobre la mejor longitud conocida, mientras que el tour de PSO es 7,2, 26,7, 84,8 y 269,0 veces más largo en $n=200$, 2\,000, 20\,000 y 200\,000.
\item \textbf{Dependencia de $n$: no hay cruce en el rango probado ($n=20$ a $n=200\,000$), y la brecha crece con $n$.} La razón $L_{PSO}/L_{ACO}$ con 2\,048 agentes y 100 it. (108 en $n=200\,000$) es 1,05 en $n=20$, 7,23 en $n=200$, 26,73 en $n=2\,000$, 84,80 en $n=20\,000$, 269,03 en $n=200\,000$: se multiplica por 256 entre $n=20$ y $n=200\,000$. No existe un tamaño de los probados en que PSO sea preferible por calidad; incluso donde la brecha es menor ($n=20$), ACO ya está en el óptimo.
\item \textbf{Tiempo: PSO es más rápido en todas las configuraciones emparejadas, pero su ventaja es menor que su desventaja de calidad.} PSO tarda entre 0,39 y 0,68 veces lo que ACO (ahorra entre 32\,\% y 61\,\% del tiempo); la razón mediana sube con $n$: 0,47 en $n=20$, 0,43 en $n=200$, 0,50 en $n=2\,000$, 0,60 en $n=20\,000$ y 0,63 en $n=200\,000$ (108 it.). En esta última instancia ACO usa 961,7 s para llegar a 385,5 y PSO 610,1 s para llegar a 103\,706: los 351,6 s que ahorra PSO no compensan un tour 269 veces más largo.
\item \textbf{Memoria: PSO usa más en todo el rango y la diferencia crece con $n$.} La razón $M_{PSO}/M_{ACO}$ con 2\,048 agentes es 1,1 en $n=20$, 2,0 en $n=200$, 9,2 en $n=2\,000$, 33,3 en $n=20\,000$, 46,2 en $n=200\,000$. En $n=200\,000$ PSO ocupa 4,7 GB y ACO 0,10 GB. La memoria de PSO sigue $3\,n\,P\cdot4$ bytes (tres vectores por partícula), por lo que con $P=20\,000$ partículas en $n=200\,000$ requeriría $\approx$48 GB y no cabe en el equipo (15,88 GB); ACO no guarda estado por hormiga entre iteraciones y su memoria casi no depende de los agentes.
\item \textbf{A igual tiempo ACO también gana.} Con el tiempo máximo común de cada $n$ ($T$), la mejor configuración de ACO y la de PSO alcanzan $L/L_{ref}$: $n=20$ ($T=0,07$ s): ACO 1,000, PSO 1,050; $n=200$ ($T=0,32$ s): ACO 1,009, PSO 7,236; $n=2\,000$ ($T=3,53$ s): ACO 1,010, PSO 26,806; $n=20\,000$ ($T=42,38$ s): ACO 1,001, PSO 84,844; $n=200\,000$ ($T=610,08$ s): ACO 1,002, PSO 269,031. Además, en 31 de 39 configuraciones de PSO existe una de ACO con tiempo y longitud menores o iguales.
\item \textbf{Veredicto.} ACO es la mejor opción para este TSP en calidad y en memoria en todo el rango probado ($n=20$ a $n=200\,000$); PSO solo gana en tiempo (0,39--0,68 veces el de ACO), una ventaja que no cambia el resultado al comparar a igual tiempo. La diferencia no depende del tamaño en el sentido del ganador, pero sí en su magnitud: es pequeña en $n=20$ y de más de dos órdenes de magnitud en $n=200\,000$.
\item \emph{Límites:} PSO con parámetros fijos ($w=0{,}7$, $c_1=c_2=1{,}5$) y sin búsqueda local; los resultados valen para esta implementación.
\end{itemize}
